% =====================================================================
%  Construction of the Simulation Environment
%  A SECTION of the FYP report. Include from the parent chapter with:
%      % =====================================================================
%  Construction of the Simulation Environment
%  A SECTION of the FYP report. Include from the parent chapter with:
%      % =====================================================================
%  Construction of the Simulation Environment
%  A SECTION of the FYP report. Include from the parent chapter with:
%      % =====================================================================
%  Construction of the Simulation Environment
%  A SECTION of the FYP report. Include from the parent chapter with:
%      \input{simulation_environment}
%
%  SCOPE. This describes the georeferenced campus world built for the
%  project. It is JOINT work: the terrain is the author's contribution,
%  the building models and the GPS integration are a project partner's.
%  Attribution is stated explicitly in the opening and must not be
%  removed -- see also the note in Section \ref{sec:world-scope} about
%  which world the HITL measurements were actually taken in.
%
%  REQUIRED PACKAGES
%      \usepackage{booktabs}
%      \usepackage{graphicx}
%      \graphicspath{{figures/}}
%
%  IMAGES EXPECTED in docs/report/figures/
%      world_rendered.png       (see checklist at end)
%      terrain_wireframe.png    (optional but recommended)
% =====================================================================

\section{Construction of the Simulation Environment}
\label{sec:world}

A georeferenced three-dimensional model of the Faculty of Engineering,
University of Ruhuna, at Hapugala, Galle, was built as the environment for
UAV simulation work. A synthetic environment would not have been adequate:
the aircraft's GPS positions, the height of the camera above the ground
beneath it, and the obstruction of radio paths by buildings all depend on the
site being real and correctly scaled.

The environment is joint work within the project. The terrain surface ---
elevation data acquisition, mesh generation and satellite texturing --- is
the present author's contribution and is described in
Section~\ref{sec:world-terrain}. The building models and the georeferenced
integration of buildings with terrain were carried out by a project partner
and are summarised in Sections~\ref{sec:world-buildings}
and~\ref{sec:world-integration} for completeness, since the environment
cannot be described coherently without them.

Every tool used is free and open source, which was a deliberate constraint:
the workflow can be reproduced for any site without licensed software.

\begin{table}[htbp]
\centering
\caption{Software used to construct the environment.}
\label{tab:world-tools}
\begin{tabular}{@{}lll@{}}
\toprule
Tool & Version & Role \\
\midrule
QGIS            & 3.22.4   & Footprint tracing, DEM clipping, imagery export \\
GDAL            & 3.4.1    & DEM reprojection and elevation sampling \\
Blender         & 5.1.1    & Mesh construction and alignment \\
BlenderGIS      & 2.2.14   & SRTM import and terrain mesh generation \\
Gazebo Classic  & 11.10.2  & Simulation environment \\
Python          & 3.10     & Coordinate conversion and validation \\
\bottomrule
\end{tabular}
\end{table}

% ---------------------------------------------------------------------
\subsection{Site and georeferencing}
\label{sec:world-site}

The world origin is placed at the front entrance of the Administration
Building, at $6.0792673^{\circ}$\,N, $80.1921607^{\circ}$\,E, approximately
$31$\,m above sea level. The surface model is \texttt{EARTH\_WGS84}, and
metre-accurate processing was carried out in UTM Zone~44N (EPSG:32644).

Anchoring the world in this way means a position in the simulation
corresponds to a position on the real campus. For a survey aircraft this is
not a convenience but a requirement: waypoints are given as GPS coordinates,
and the autopilot's own position estimate is a GPS solution, so a world whose
origin is arbitrary cannot be flown with real mission data.

\begin{table}[htbp]
\centering
\caption{Environmental parameters, set for Galle rather than left at default.}
\label{tab:world-params}
\begin{tabular}{@{}llp{58mm}@{}}
\toprule
Parameter & Value & Why it is set this way \\
\midrule
Gravity
  & $-9.7804$\,m\,s$^{-2}$
  & Local value; the $-9.80$ default overstates weight near the equator \\[3pt]
Magnetic field (WMM2025)
  & $(6.0,\ 23.0,\ -38.0)\ \mu$T
  & The autopilot derives heading from the magnetometer; a wrong field
    produces a persistent heading bias \\[3pt]
Temperature
  & $301$\,K ($28^{\circ}$C)
  & Tropical, sets air density and therefore rotor thrust \\[3pt]
Pressure
  & $101{,}000$\,Pa
  & At the $31$\,m datum \\[3pt]
Lapse rate
  & $-0.0065$\,K\,m$^{-1}$
  & Density falls with altitude as the aircraft climbs \\[3pt]
Wind
  & $0$ (baseline)
  & Zero for repeatable measurement; the south-west monsoon condition
    is available as a separate case \\
\bottomrule
\end{tabular}
\end{table}

Wind is held at zero for the measurements reported in this work. A moving
air mass makes each flight follow a slightly different path, and the
experiment compares detection performance between configurations rather than
robustness to disturbance, so a still atmosphere removes a source of
run-to-run variation. A south-west monsoon case, appropriate to Galle between
May and September, is configured separately for later use.

% ---------------------------------------------------------------------
\subsection{Terrain}
\label{sec:world-terrain}

Elevation data was taken from the Shuttle Radar Topography Mission, the
NASA-led survey that produced near-global digital elevation models at
$30$\,m horizontal resolution. SRTM was chosen over the alternatives ---
Copernicus Open Access Hub and USGS Earth Explorer among them --- because it
gives complete coverage of the area with no missing tiles. Cloud cover is
the usual reason such datasets have gaps, since the sensors cannot measure
terrain through it, and a dataset must be inspected for missing data before
it is used rather than after.

The DEM was clipped in QGIS from an $8 \times 8$\,km tile to the campus
extent and reprojected from EPSG:4326 to EPSG:32644 so that subsequent
processing was in metres rather than degrees. Elevation across the clipped
area ranges from $6$\,m to $52$\,m above sea level.

The clipped DEM was imported into Blender through the BlenderGIS add-on,
which reads elevation products directly and generates a triangulated surface
from them. The delivered mesh is a regular $65 \times 65$ grid of $4{,}225$
vertices covering $382$\,m east--west by $535$\,m north--south, with
$42.7$\,m of relief across it.

Two properties of this surface deserve comment.

\paragraph{Resolution.}
The grid spacing of roughly $6$\,m is finer than SRTM's native $30$\,m
because BlenderGIS resamples the elevation product onto its own grid. The
mesh is therefore smoother than the source data, not more accurate than it:
SRTM's specified vertical accuracy is $\pm 16$\,m, and no processing step
improves on that. What the terrain represents faithfully is the \emph{shape}
of the site --- where the ground rises and falls --- rather than absolute
height at any single point. That is what the experiment requires, since the
aircraft flies at a commanded height above its launch point and it is the
variation in ground clearance across a patrol that changes the apparent size
of a person in the image.

\paragraph{Export format.}
The terrain was exported in COLLADA rather than STL. The distinction
resolves two difficulties that arise with STL in this workflow. First, STL
carries no unit declaration, so a mesh must be rescaled by hand after import
and an error there propagates into every distance the simulation reports;
the COLLADA file declares \texttt{<unit name="meter" meter="1"/>} with a
$Z$-up axis, so Gazebo reads it at true scale unmodified. Second, STL stores
only geometry and has no texture coordinates, so an STL terrain cannot carry
aerial imagery and renders as a bare surface. COLLADA stores per-vertex
texture coordinates alongside the geometry, which is what allows the imagery
described next to be applied at all.

\paragraph{Imagery.}
Aerial imagery was exported from Google Satellite through QGIS XYZ tile
layers in Web Mercator projection, at $1792 \times 1280$ pixels, and applied
to the terrain as a texture. A cloud-free export was selected: cloud in the
imagery would appear as permanent white patches on the ground and would be a
standing source of false detections. Beyond appearance, the texture supplies
the visual clutter of vegetation, paths and roofs against which the detector
must distinguish a person --- a uniformly coloured surface would make the
detection task easier than it is in reality, and any recall figure measured
against it would be optimistic.

The imagery was initially referenced in GeoTIFF form, as exported. That file
carries EXIF metadata which Gazebo's image library does not recognise,
producing roughly thirty diagnostic messages on every load; the same image in
PNG at identical resolution was substituted, removing them at no cost in
quality. The messages were harmless, but indistinguishable at a glance from
a genuine loading failure.

% ---------------------------------------------------------------------
\subsection{Structures}
\label{sec:world-buildings}

\emph{The work described in this subsection and the next was carried out by a
project partner. It is included because the environment cannot be described
without it.}

Seventeen building footprints were traced manually in QGIS over a Google
Satellite tile layer, each carrying attributes for name, height, floor count
and roof type, and exported as GeoJSON in WGS84. The footprints were imported
into Blender, extruded to height, and given terracotta hip roofs modelled
individually.

Manual tracing was necessary because building outlines for this campus are
not available at usable quality in open mapping databases. The cost of that
choice is that footprint accuracy is limited by how precisely a polygon can
be placed over satellite imagery by hand, which is one of the two
contributions to the position error reported below.

% ---------------------------------------------------------------------
\subsection{Integration and positional accuracy}
\label{sec:world-integration}

Terrain and buildings were produced independently and against different
origins --- the terrain against $6.07778^{\circ}$\,N,
$80.20722^{\circ}$\,E and the buildings against the Administration Building
entrance --- so an offset of $1666.281$\,m east and $-166.891$\,m north was
computed and applied to bring them into one coordinate system. Each building
was then adjusted vertically to sit on the terrain slope beneath it, and the
combined scene exported as a single mesh. The model is placed at
$z = -31.412$\,m so that the terrain surface coincides with the simulator's
$z = 0$ plane.

A $2000 \times 2000$\,m invisible collision plane is placed at that same
level. Without it, an aircraft leaving the $382 \times 535$\,m modelled area
would encounter no ground at all and fall indefinitely, ending the run.

Positional accuracy was checked against two references. Against the source
GeoJSON footprints, agreement was below $0.01$\,m, confirming that the
coordinate transformation itself is correct. Against building centroids
picked independently from Google Maps, the mean error over sixteen structures
was $5.64$\,m, with a maximum of $10.59$\,m and a minimum of $0.38$\,m.

The second figure is the meaningful one, and it is consistent with what the
inputs allow: SRTM's $30$\,m cells and the precision of manual tracing
together account for errors of this order. For context, consumer GPS
receivers of the kind carried by small UAVs achieve $2$--$5$\,m in practice,
so the discrepancy between the modelled campus and the real one is of the
same order as the uncertainty a real aircraft would have in its own position.

% ---------------------------------------------------------------------
\subsection{Relationship to the experiments reported in this work}
\label{sec:world-scope}

The campus environment described above is a contribution in its own right and
is intended for mission-level UAV work at this site. The hardware-in-the-loop
measurements reported in Section~\ref{sec:hitl} were \emph{not} taken in it.
They were taken in a separate world, which was already validated against the
network-namespace and lockstep configuration those experiments depend on, and
which contains five human models at fixed coordinates that make detection
recall a counted quantity.

The two are kept distinct deliberately. Changing the world and the network
architecture at the same time would have made any change in the results
impossible to attribute. Running the hardware-in-the-loop experiments in the
campus environment is the natural next step, and requires only that the
autopilot's flight-dynamics plugin be added to that world.

\begin{figure}[htbp]
\centering
\includegraphics[width=0.9\textwidth]{world_rendered.png}
\caption{The campus environment in Gazebo. Google Satellite imagery is
applied to an SRTM-derived terrain mesh, with individually modelled buildings
placed on it. Terrain and structures both carry collision geometry generated
from the meshes used for rendering, so the surface an aircraft can strike is
the surface its camera sees.}
\label{fig:world-rendered}
\end{figure}

% =====================================================================
%  IMAGE CHECKLIST -- save into docs/report/figures/
%
%  world_rendered.png   (required)
%      Gazebo client view of the campus. Two things to fix relative to the
%      screenshot taken so far:
%        1. Pull the camera back and up. The current shot is close in on a
%           single roof; a wider view showing several buildings and the
%           terrain relief demonstrates far more.
%        2. Crop out the left World/Insert/Layers panel and the bottom
%           status bar. Keep only the 3D viewport.
%
%  terrain_wireframe.png   (optional, recommended)
%      View > Wireframe in Gazebo, camera angled low across the terrain.
%      This makes the 65x65 grid and the relief visible at once, and is the
%      only image that proves the terrain is a height map rather than a
%      textured plane -- the claim in this section a reader is most likely
%      to doubt.
%
%  NOT recommended as a figure: the GPS validation table. It is a partner's
%  result; cite the summary figures in prose as done above rather than
%  reproducing the full table under your own name.
% =====================================================================

%
%  SCOPE. This describes the georeferenced campus world built for the
%  project. It is JOINT work: the terrain is the author's contribution,
%  the building models and the GPS integration are a project partner's.
%  Attribution is stated explicitly in the opening and must not be
%  removed -- see also the note in Section \ref{sec:world-scope} about
%  which world the HITL measurements were actually taken in.
%
%  REQUIRED PACKAGES
%      \usepackage{booktabs}
%      \usepackage{graphicx}
%      \graphicspath{{figures/}}
%
%  IMAGES EXPECTED in docs/report/figures/
%      world_rendered.png       (see checklist at end)
%      terrain_wireframe.png    (optional but recommended)
% =====================================================================

\section{Construction of the Simulation Environment}
\label{sec:world}

A georeferenced three-dimensional model of the Faculty of Engineering,
University of Ruhuna, at Hapugala, Galle, was built as the environment for
UAV simulation work. A synthetic environment would not have been adequate:
the aircraft's GPS positions, the height of the camera above the ground
beneath it, and the obstruction of radio paths by buildings all depend on the
site being real and correctly scaled.

The environment is joint work within the project. The terrain surface ---
elevation data acquisition, mesh generation and satellite texturing --- is
the present author's contribution and is described in
Section~\ref{sec:world-terrain}. The building models and the georeferenced
integration of buildings with terrain were carried out by a project partner
and are summarised in Sections~\ref{sec:world-buildings}
and~\ref{sec:world-integration} for completeness, since the environment
cannot be described coherently without them.

Every tool used is free and open source, which was a deliberate constraint:
the workflow can be reproduced for any site without licensed software.

\begin{table}[htbp]
\centering
\caption{Software used to construct the environment.}
\label{tab:world-tools}
\begin{tabular}{@{}lll@{}}
\toprule
Tool & Version & Role \\
\midrule
QGIS            & 3.22.4   & Footprint tracing, DEM clipping, imagery export \\
GDAL            & 3.4.1    & DEM reprojection and elevation sampling \\
Blender         & 5.1.1    & Mesh construction and alignment \\
BlenderGIS      & 2.2.14   & SRTM import and terrain mesh generation \\
Gazebo Classic  & 11.10.2  & Simulation environment \\
Python          & 3.10     & Coordinate conversion and validation \\
\bottomrule
\end{tabular}
\end{table}

% ---------------------------------------------------------------------
\subsection{Site and georeferencing}
\label{sec:world-site}

The world origin is placed at the front entrance of the Administration
Building, at $6.0792673^{\circ}$\,N, $80.1921607^{\circ}$\,E, approximately
$31$\,m above sea level. The surface model is \texttt{EARTH\_WGS84}, and
metre-accurate processing was carried out in UTM Zone~44N (EPSG:32644).

Anchoring the world in this way means a position in the simulation
corresponds to a position on the real campus. For a survey aircraft this is
not a convenience but a requirement: waypoints are given as GPS coordinates,
and the autopilot's own position estimate is a GPS solution, so a world whose
origin is arbitrary cannot be flown with real mission data.

\begin{table}[htbp]
\centering
\caption{Environmental parameters, set for Galle rather than left at default.}
\label{tab:world-params}
\begin{tabular}{@{}llp{58mm}@{}}
\toprule
Parameter & Value & Why it is set this way \\
\midrule
Gravity
  & $-9.7804$\,m\,s$^{-2}$
  & Local value; the $-9.80$ default overstates weight near the equator \\[3pt]
Magnetic field (WMM2025)
  & $(6.0,\ 23.0,\ -38.0)\ \mu$T
  & The autopilot derives heading from the magnetometer; a wrong field
    produces a persistent heading bias \\[3pt]
Temperature
  & $301$\,K ($28^{\circ}$C)
  & Tropical, sets air density and therefore rotor thrust \\[3pt]
Pressure
  & $101{,}000$\,Pa
  & At the $31$\,m datum \\[3pt]
Lapse rate
  & $-0.0065$\,K\,m$^{-1}$
  & Density falls with altitude as the aircraft climbs \\[3pt]
Wind
  & $0$ (baseline)
  & Zero for repeatable measurement; the south-west monsoon condition
    is available as a separate case \\
\bottomrule
\end{tabular}
\end{table}

Wind is held at zero for the measurements reported in this work. A moving
air mass makes each flight follow a slightly different path, and the
experiment compares detection performance between configurations rather than
robustness to disturbance, so a still atmosphere removes a source of
run-to-run variation. A south-west monsoon case, appropriate to Galle between
May and September, is configured separately for later use.

% ---------------------------------------------------------------------
\subsection{Terrain}
\label{sec:world-terrain}

Elevation data was taken from the Shuttle Radar Topography Mission, the
NASA-led survey that produced near-global digital elevation models at
$30$\,m horizontal resolution. SRTM was chosen over the alternatives ---
Copernicus Open Access Hub and USGS Earth Explorer among them --- because it
gives complete coverage of the area with no missing tiles. Cloud cover is
the usual reason such datasets have gaps, since the sensors cannot measure
terrain through it, and a dataset must be inspected for missing data before
it is used rather than after.

The DEM was clipped in QGIS from an $8 \times 8$\,km tile to the campus
extent and reprojected from EPSG:4326 to EPSG:32644 so that subsequent
processing was in metres rather than degrees. Elevation across the clipped
area ranges from $6$\,m to $52$\,m above sea level.

The clipped DEM was imported into Blender through the BlenderGIS add-on,
which reads elevation products directly and generates a triangulated surface
from them. The delivered mesh is a regular $65 \times 65$ grid of $4{,}225$
vertices covering $382$\,m east--west by $535$\,m north--south, with
$42.7$\,m of relief across it.

Two properties of this surface deserve comment.

\paragraph{Resolution.}
The grid spacing of roughly $6$\,m is finer than SRTM's native $30$\,m
because BlenderGIS resamples the elevation product onto its own grid. The
mesh is therefore smoother than the source data, not more accurate than it:
SRTM's specified vertical accuracy is $\pm 16$\,m, and no processing step
improves on that. What the terrain represents faithfully is the \emph{shape}
of the site --- where the ground rises and falls --- rather than absolute
height at any single point. That is what the experiment requires, since the
aircraft flies at a commanded height above its launch point and it is the
variation in ground clearance across a patrol that changes the apparent size
of a person in the image.

\paragraph{Export format.}
The terrain was exported in COLLADA rather than STL. The distinction
resolves two difficulties that arise with STL in this workflow. First, STL
carries no unit declaration, so a mesh must be rescaled by hand after import
and an error there propagates into every distance the simulation reports;
the COLLADA file declares \texttt{<unit name="meter" meter="1"/>} with a
$Z$-up axis, so Gazebo reads it at true scale unmodified. Second, STL stores
only geometry and has no texture coordinates, so an STL terrain cannot carry
aerial imagery and renders as a bare surface. COLLADA stores per-vertex
texture coordinates alongside the geometry, which is what allows the imagery
described next to be applied at all.

\paragraph{Imagery.}
Aerial imagery was exported from Google Satellite through QGIS XYZ tile
layers in Web Mercator projection, at $1792 \times 1280$ pixels, and applied
to the terrain as a texture. A cloud-free export was selected: cloud in the
imagery would appear as permanent white patches on the ground and would be a
standing source of false detections. Beyond appearance, the texture supplies
the visual clutter of vegetation, paths and roofs against which the detector
must distinguish a person --- a uniformly coloured surface would make the
detection task easier than it is in reality, and any recall figure measured
against it would be optimistic.

The imagery was initially referenced in GeoTIFF form, as exported. That file
carries EXIF metadata which Gazebo's image library does not recognise,
producing roughly thirty diagnostic messages on every load; the same image in
PNG at identical resolution was substituted, removing them at no cost in
quality. The messages were harmless, but indistinguishable at a glance from
a genuine loading failure.

% ---------------------------------------------------------------------
\subsection{Structures}
\label{sec:world-buildings}

\emph{The work described in this subsection and the next was carried out by a
project partner. It is included because the environment cannot be described
without it.}

Seventeen building footprints were traced manually in QGIS over a Google
Satellite tile layer, each carrying attributes for name, height, floor count
and roof type, and exported as GeoJSON in WGS84. The footprints were imported
into Blender, extruded to height, and given terracotta hip roofs modelled
individually.

Manual tracing was necessary because building outlines for this campus are
not available at usable quality in open mapping databases. The cost of that
choice is that footprint accuracy is limited by how precisely a polygon can
be placed over satellite imagery by hand, which is one of the two
contributions to the position error reported below.

% ---------------------------------------------------------------------
\subsection{Integration and positional accuracy}
\label{sec:world-integration}

Terrain and buildings were produced independently and against different
origins --- the terrain against $6.07778^{\circ}$\,N,
$80.20722^{\circ}$\,E and the buildings against the Administration Building
entrance --- so an offset of $1666.281$\,m east and $-166.891$\,m north was
computed and applied to bring them into one coordinate system. Each building
was then adjusted vertically to sit on the terrain slope beneath it, and the
combined scene exported as a single mesh. The model is placed at
$z = -31.412$\,m so that the terrain surface coincides with the simulator's
$z = 0$ plane.

A $2000 \times 2000$\,m invisible collision plane is placed at that same
level. Without it, an aircraft leaving the $382 \times 535$\,m modelled area
would encounter no ground at all and fall indefinitely, ending the run.

Positional accuracy was checked against two references. Against the source
GeoJSON footprints, agreement was below $0.01$\,m, confirming that the
coordinate transformation itself is correct. Against building centroids
picked independently from Google Maps, the mean error over sixteen structures
was $5.64$\,m, with a maximum of $10.59$\,m and a minimum of $0.38$\,m.

The second figure is the meaningful one, and it is consistent with what the
inputs allow: SRTM's $30$\,m cells and the precision of manual tracing
together account for errors of this order. For context, consumer GPS
receivers of the kind carried by small UAVs achieve $2$--$5$\,m in practice,
so the discrepancy between the modelled campus and the real one is of the
same order as the uncertainty a real aircraft would have in its own position.

% ---------------------------------------------------------------------
\subsection{Relationship to the experiments reported in this work}
\label{sec:world-scope}

The campus environment described above is a contribution in its own right and
is intended for mission-level UAV work at this site. The hardware-in-the-loop
measurements reported in Section~\ref{sec:hitl} were \emph{not} taken in it.
They were taken in a separate world, which was already validated against the
network-namespace and lockstep configuration those experiments depend on, and
which contains five human models at fixed coordinates that make detection
recall a counted quantity.

The two are kept distinct deliberately. Changing the world and the network
architecture at the same time would have made any change in the results
impossible to attribute. Running the hardware-in-the-loop experiments in the
campus environment is the natural next step, and requires only that the
autopilot's flight-dynamics plugin be added to that world.

\begin{figure}[htbp]
\centering
\includegraphics[width=0.9\textwidth]{world_rendered.png}
\caption{The campus environment in Gazebo. Google Satellite imagery is
applied to an SRTM-derived terrain mesh, with individually modelled buildings
placed on it. Terrain and structures both carry collision geometry generated
from the meshes used for rendering, so the surface an aircraft can strike is
the surface its camera sees.}
\label{fig:world-rendered}
\end{figure}

% =====================================================================
%  IMAGE CHECKLIST -- save into docs/report/figures/
%
%  world_rendered.png   (required)
%      Gazebo client view of the campus. Two things to fix relative to the
%      screenshot taken so far:
%        1. Pull the camera back and up. The current shot is close in on a
%           single roof; a wider view showing several buildings and the
%           terrain relief demonstrates far more.
%        2. Crop out the left World/Insert/Layers panel and the bottom
%           status bar. Keep only the 3D viewport.
%
%  terrain_wireframe.png   (optional, recommended)
%      View > Wireframe in Gazebo, camera angled low across the terrain.
%      This makes the 65x65 grid and the relief visible at once, and is the
%      only image that proves the terrain is a height map rather than a
%      textured plane -- the claim in this section a reader is most likely
%      to doubt.
%
%  NOT recommended as a figure: the GPS validation table. It is a partner's
%  result; cite the summary figures in prose as done above rather than
%  reproducing the full table under your own name.
% =====================================================================

%
%  SCOPE. This describes the georeferenced campus world built for the
%  project. It is JOINT work: the terrain is the author's contribution,
%  the building models and the GPS integration are a project partner's.
%  Attribution is stated explicitly in the opening and must not be
%  removed -- see also the note in Section \ref{sec:world-scope} about
%  which world the HITL measurements were actually taken in.
%
%  REQUIRED PACKAGES
%      \usepackage{booktabs}
%      \usepackage{graphicx}
%      \graphicspath{{figures/}}
%
%  IMAGES EXPECTED in docs/report/figures/
%      world_rendered.png       (see checklist at end)
%      terrain_wireframe.png    (optional but recommended)
% =====================================================================

\section{Construction of the Simulation Environment}
\label{sec:world}

A georeferenced three-dimensional model of the Faculty of Engineering,
University of Ruhuna, at Hapugala, Galle, was built as the environment for
UAV simulation work. A synthetic environment would not have been adequate:
the aircraft's GPS positions, the height of the camera above the ground
beneath it, and the obstruction of radio paths by buildings all depend on the
site being real and correctly scaled.

The environment is joint work within the project. The terrain surface ---
elevation data acquisition, mesh generation and satellite texturing --- is
the present author's contribution and is described in
Section~\ref{sec:world-terrain}. The building models and the georeferenced
integration of buildings with terrain were carried out by a project partner
and are summarised in Sections~\ref{sec:world-buildings}
and~\ref{sec:world-integration} for completeness, since the environment
cannot be described coherently without them.

Every tool used is free and open source, which was a deliberate constraint:
the workflow can be reproduced for any site without licensed software.

\begin{table}[htbp]
\centering
\caption{Software used to construct the environment.}
\label{tab:world-tools}
\begin{tabular}{@{}lll@{}}
\toprule
Tool & Version & Role \\
\midrule
QGIS            & 3.22.4   & Footprint tracing, DEM clipping, imagery export \\
GDAL            & 3.4.1    & DEM reprojection and elevation sampling \\
Blender         & 5.1.1    & Mesh construction and alignment \\
BlenderGIS      & 2.2.14   & SRTM import and terrain mesh generation \\
Gazebo Classic  & 11.10.2  & Simulation environment \\
Python          & 3.10     & Coordinate conversion and validation \\
\bottomrule
\end{tabular}
\end{table}

% ---------------------------------------------------------------------
\subsection{Site and georeferencing}
\label{sec:world-site}

The world origin is placed at the front entrance of the Administration
Building, at $6.0792673^{\circ}$\,N, $80.1921607^{\circ}$\,E, approximately
$31$\,m above sea level. The surface model is \texttt{EARTH\_WGS84}, and
metre-accurate processing was carried out in UTM Zone~44N (EPSG:32644).

Anchoring the world in this way means a position in the simulation
corresponds to a position on the real campus. For a survey aircraft this is
not a convenience but a requirement: waypoints are given as GPS coordinates,
and the autopilot's own position estimate is a GPS solution, so a world whose
origin is arbitrary cannot be flown with real mission data.

\begin{table}[htbp]
\centering
\caption{Environmental parameters, set for Galle rather than left at default.}
\label{tab:world-params}
\begin{tabular}{@{}llp{58mm}@{}}
\toprule
Parameter & Value & Why it is set this way \\
\midrule
Gravity
  & $-9.7804$\,m\,s$^{-2}$
  & Local value; the $-9.80$ default overstates weight near the equator \\[3pt]
Magnetic field (WMM2025)
  & $(6.0,\ 23.0,\ -38.0)\ \mu$T
  & The autopilot derives heading from the magnetometer; a wrong field
    produces a persistent heading bias \\[3pt]
Temperature
  & $301$\,K ($28^{\circ}$C)
  & Tropical, sets air density and therefore rotor thrust \\[3pt]
Pressure
  & $101{,}000$\,Pa
  & At the $31$\,m datum \\[3pt]
Lapse rate
  & $-0.0065$\,K\,m$^{-1}$
  & Density falls with altitude as the aircraft climbs \\[3pt]
Wind
  & $0$ (baseline)
  & Zero for repeatable measurement; the south-west monsoon condition
    is available as a separate case \\
\bottomrule
\end{tabular}
\end{table}

Wind is held at zero for the measurements reported in this work. A moving
air mass makes each flight follow a slightly different path, and the
experiment compares detection performance between configurations rather than
robustness to disturbance, so a still atmosphere removes a source of
run-to-run variation. A south-west monsoon case, appropriate to Galle between
May and September, is configured separately for later use.

% ---------------------------------------------------------------------
\subsection{Terrain}
\label{sec:world-terrain}

Elevation data was taken from the Shuttle Radar Topography Mission, the
NASA-led survey that produced near-global digital elevation models at
$30$\,m horizontal resolution. SRTM was chosen over the alternatives ---
Copernicus Open Access Hub and USGS Earth Explorer among them --- because it
gives complete coverage of the area with no missing tiles. Cloud cover is
the usual reason such datasets have gaps, since the sensors cannot measure
terrain through it, and a dataset must be inspected for missing data before
it is used rather than after.

The DEM was clipped in QGIS from an $8 \times 8$\,km tile to the campus
extent and reprojected from EPSG:4326 to EPSG:32644 so that subsequent
processing was in metres rather than degrees. Elevation across the clipped
area ranges from $6$\,m to $52$\,m above sea level.

The clipped DEM was imported into Blender through the BlenderGIS add-on,
which reads elevation products directly and generates a triangulated surface
from them. The delivered mesh is a regular $65 \times 65$ grid of $4{,}225$
vertices covering $382$\,m east--west by $535$\,m north--south, with
$42.7$\,m of relief across it.

Two properties of this surface deserve comment.

\paragraph{Resolution.}
The grid spacing of roughly $6$\,m is finer than SRTM's native $30$\,m
because BlenderGIS resamples the elevation product onto its own grid. The
mesh is therefore smoother than the source data, not more accurate than it:
SRTM's specified vertical accuracy is $\pm 16$\,m, and no processing step
improves on that. What the terrain represents faithfully is the \emph{shape}
of the site --- where the ground rises and falls --- rather than absolute
height at any single point. That is what the experiment requires, since the
aircraft flies at a commanded height above its launch point and it is the
variation in ground clearance across a patrol that changes the apparent size
of a person in the image.

\paragraph{Export format.}
The terrain was exported in COLLADA rather than STL. The distinction
resolves two difficulties that arise with STL in this workflow. First, STL
carries no unit declaration, so a mesh must be rescaled by hand after import
and an error there propagates into every distance the simulation reports;
the COLLADA file declares \texttt{<unit name="meter" meter="1"/>} with a
$Z$-up axis, so Gazebo reads it at true scale unmodified. Second, STL stores
only geometry and has no texture coordinates, so an STL terrain cannot carry
aerial imagery and renders as a bare surface. COLLADA stores per-vertex
texture coordinates alongside the geometry, which is what allows the imagery
described next to be applied at all.

\paragraph{Imagery.}
Aerial imagery was exported from Google Satellite through QGIS XYZ tile
layers in Web Mercator projection, at $1792 \times 1280$ pixels, and applied
to the terrain as a texture. A cloud-free export was selected: cloud in the
imagery would appear as permanent white patches on the ground and would be a
standing source of false detections. Beyond appearance, the texture supplies
the visual clutter of vegetation, paths and roofs against which the detector
must distinguish a person --- a uniformly coloured surface would make the
detection task easier than it is in reality, and any recall figure measured
against it would be optimistic.

The imagery was initially referenced in GeoTIFF form, as exported. That file
carries EXIF metadata which Gazebo's image library does not recognise,
producing roughly thirty diagnostic messages on every load; the same image in
PNG at identical resolution was substituted, removing them at no cost in
quality. The messages were harmless, but indistinguishable at a glance from
a genuine loading failure.

% ---------------------------------------------------------------------
\subsection{Structures}
\label{sec:world-buildings}

\emph{The work described in this subsection and the next was carried out by a
project partner. It is included because the environment cannot be described
without it.}

Seventeen building footprints were traced manually in QGIS over a Google
Satellite tile layer, each carrying attributes for name, height, floor count
and roof type, and exported as GeoJSON in WGS84. The footprints were imported
into Blender, extruded to height, and given terracotta hip roofs modelled
individually.

Manual tracing was necessary because building outlines for this campus are
not available at usable quality in open mapping databases. The cost of that
choice is that footprint accuracy is limited by how precisely a polygon can
be placed over satellite imagery by hand, which is one of the two
contributions to the position error reported below.

% ---------------------------------------------------------------------
\subsection{Integration and positional accuracy}
\label{sec:world-integration}

Terrain and buildings were produced independently and against different
origins --- the terrain against $6.07778^{\circ}$\,N,
$80.20722^{\circ}$\,E and the buildings against the Administration Building
entrance --- so an offset of $1666.281$\,m east and $-166.891$\,m north was
computed and applied to bring them into one coordinate system. Each building
was then adjusted vertically to sit on the terrain slope beneath it, and the
combined scene exported as a single mesh. The model is placed at
$z = -31.412$\,m so that the terrain surface coincides with the simulator's
$z = 0$ plane.

A $2000 \times 2000$\,m invisible collision plane is placed at that same
level. Without it, an aircraft leaving the $382 \times 535$\,m modelled area
would encounter no ground at all and fall indefinitely, ending the run.

Positional accuracy was checked against two references. Against the source
GeoJSON footprints, agreement was below $0.01$\,m, confirming that the
coordinate transformation itself is correct. Against building centroids
picked independently from Google Maps, the mean error over sixteen structures
was $5.64$\,m, with a maximum of $10.59$\,m and a minimum of $0.38$\,m.

The second figure is the meaningful one, and it is consistent with what the
inputs allow: SRTM's $30$\,m cells and the precision of manual tracing
together account for errors of this order. For context, consumer GPS
receivers of the kind carried by small UAVs achieve $2$--$5$\,m in practice,
so the discrepancy between the modelled campus and the real one is of the
same order as the uncertainty a real aircraft would have in its own position.

% ---------------------------------------------------------------------
\subsection{Relationship to the experiments reported in this work}
\label{sec:world-scope}

The campus environment described above is a contribution in its own right and
is intended for mission-level UAV work at this site. The hardware-in-the-loop
measurements reported in Section~\ref{sec:hitl} were \emph{not} taken in it.
They were taken in a separate world, which was already validated against the
network-namespace and lockstep configuration those experiments depend on, and
which contains five human models at fixed coordinates that make detection
recall a counted quantity.

The two are kept distinct deliberately. Changing the world and the network
architecture at the same time would have made any change in the results
impossible to attribute. Running the hardware-in-the-loop experiments in the
campus environment is the natural next step, and requires only that the
autopilot's flight-dynamics plugin be added to that world.

\begin{figure}[htbp]
\centering
\includegraphics[width=0.9\textwidth]{world_rendered.png}
\caption{The campus environment in Gazebo. Google Satellite imagery is
applied to an SRTM-derived terrain mesh, with individually modelled buildings
placed on it. Terrain and structures both carry collision geometry generated
from the meshes used for rendering, so the surface an aircraft can strike is
the surface its camera sees.}
\label{fig:world-rendered}
\end{figure}

% =====================================================================
%  IMAGE CHECKLIST -- save into docs/report/figures/
%
%  world_rendered.png   (required)
%      Gazebo client view of the campus. Two things to fix relative to the
%      screenshot taken so far:
%        1. Pull the camera back and up. The current shot is close in on a
%           single roof; a wider view showing several buildings and the
%           terrain relief demonstrates far more.
%        2. Crop out the left World/Insert/Layers panel and the bottom
%           status bar. Keep only the 3D viewport.
%
%  terrain_wireframe.png   (optional, recommended)
%      View > Wireframe in Gazebo, camera angled low across the terrain.
%      This makes the 65x65 grid and the relief visible at once, and is the
%      only image that proves the terrain is a height map rather than a
%      textured plane -- the claim in this section a reader is most likely
%      to doubt.
%
%  NOT recommended as a figure: the GPS validation table. It is a partner's
%  result; cite the summary figures in prose as done above rather than
%  reproducing the full table under your own name.
% =====================================================================

%
%  SCOPE. This describes the georeferenced campus world built for the
%  project. It is JOINT work: the terrain is the author's contribution,
%  the building models and the GPS integration are a project partner's.
%  Attribution is stated explicitly in the opening and must not be
%  removed -- see also the note in Section \ref{sec:world-scope} about
%  which world the HITL measurements were actually taken in.
%
%  REQUIRED PACKAGES
%      \usepackage{booktabs}
%      \usepackage{graphicx}
%      \graphicspath{{figures/}}
%
%  IMAGES EXPECTED in docs/report/figures/
%      world_rendered.png       (see checklist at end)
%      terrain_wireframe.png    (optional but recommended)
% =====================================================================

\section{Construction of the Simulation Environment}
\label{sec:world}

A georeferenced three-dimensional model of the Faculty of Engineering,
University of Ruhuna, at Hapugala, Galle, was built as the environment for
UAV simulation work. A synthetic environment would not have been adequate:
the aircraft's GPS positions, the height of the camera above the ground
beneath it, and the obstruction of radio paths by buildings all depend on the
site being real and correctly scaled.

The environment is joint work within the project. The terrain surface ---
elevation data acquisition, mesh generation and satellite texturing --- is
the present author's contribution and is described in
Section~\ref{sec:world-terrain}. The building models and the georeferenced
integration of buildings with terrain were carried out by a project partner
and are summarised in Sections~\ref{sec:world-buildings}
and~\ref{sec:world-integration} for completeness, since the environment
cannot be described coherently without them.

Every tool used is free and open source, which was a deliberate constraint:
the workflow can be reproduced for any site without licensed software.

\begin{table}[htbp]
\centering
\caption{Software used to construct the environment.}
\label{tab:world-tools}
\begin{tabular}{@{}lll@{}}
\toprule
Tool & Version & Role \\
\midrule
QGIS            & 3.22.4   & Footprint tracing, DEM clipping, imagery export \\
GDAL            & 3.4.1    & DEM reprojection and elevation sampling \\
Blender         & 5.1.1    & Mesh construction and alignment \\
BlenderGIS      & 2.2.14   & SRTM import and terrain mesh generation \\
Gazebo Classic  & 11.10.2  & Simulation environment \\
Python          & 3.10     & Coordinate conversion and validation \\
\bottomrule
\end{tabular}
\end{table}

% ---------------------------------------------------------------------
\subsection{Site and georeferencing}
\label{sec:world-site}

The world origin is placed at the front entrance of the Administration
Building, at $6.0792673^{\circ}$\,N, $80.1921607^{\circ}$\,E, approximately
$31$\,m above sea level. The surface model is \texttt{EARTH\_WGS84}, and
metre-accurate processing was carried out in UTM Zone~44N (EPSG:32644).

Anchoring the world in this way means a position in the simulation
corresponds to a position on the real campus. For a survey aircraft this is
not a convenience but a requirement: waypoints are given as GPS coordinates,
and the autopilot's own position estimate is a GPS solution, so a world whose
origin is arbitrary cannot be flown with real mission data.

\begin{table}[htbp]
\centering
\caption{Environmental parameters, set for Galle rather than left at default.}
\label{tab:world-params}
\begin{tabular}{@{}llp{58mm}@{}}
\toprule
Parameter & Value & Why it is set this way \\
\midrule
Gravity
  & $-9.7804$\,m\,s$^{-2}$
  & Local value; the $-9.80$ default overstates weight near the equator \\[3pt]
Magnetic field (WMM2025)
  & $(6.0,\ 23.0,\ -38.0)\ \mu$T
  & The autopilot derives heading from the magnetometer; a wrong field
    produces a persistent heading bias \\[3pt]
Temperature
  & $301$\,K ($28^{\circ}$C)
  & Tropical, sets air density and therefore rotor thrust \\[3pt]
Pressure
  & $101{,}000$\,Pa
  & At the $31$\,m datum \\[3pt]
Lapse rate
  & $-0.0065$\,K\,m$^{-1}$
  & Density falls with altitude as the aircraft climbs \\[3pt]
Wind
  & $0$ (baseline)
  & Zero for repeatable measurement; the south-west monsoon condition
    is available as a separate case \\
\bottomrule
\end{tabular}
\end{table}

Wind is held at zero for the measurements reported in this work. A moving
air mass makes each flight follow a slightly different path, and the
experiment compares detection performance between configurations rather than
robustness to disturbance, so a still atmosphere removes a source of
run-to-run variation. A south-west monsoon case, appropriate to Galle between
May and September, is configured separately for later use.

% ---------------------------------------------------------------------
\subsection{Terrain}
\label{sec:world-terrain}

Elevation data was taken from the Shuttle Radar Topography Mission, the
NASA-led survey that produced near-global digital elevation models at
$30$\,m horizontal resolution. SRTM was chosen over the alternatives ---
Copernicus Open Access Hub and USGS Earth Explorer among them --- because it
gives complete coverage of the area with no missing tiles. Cloud cover is
the usual reason such datasets have gaps, since the sensors cannot measure
terrain through it, and a dataset must be inspected for missing data before
it is used rather than after.

The DEM was clipped in QGIS from an $8 \times 8$\,km tile to the campus
extent and reprojected from EPSG:4326 to EPSG:32644 so that subsequent
processing was in metres rather than degrees. Elevation across the clipped
area ranges from $6$\,m to $52$\,m above sea level.

The clipped DEM was imported into Blender through the BlenderGIS add-on,
which reads elevation products directly and generates a triangulated surface
from them. The delivered mesh is a regular $65 \times 65$ grid of $4{,}225$
vertices covering $382$\,m east--west by $535$\,m north--south, with
$42.7$\,m of relief across it.

Two properties of this surface deserve comment.

\paragraph{Resolution.}
The grid spacing of roughly $6$\,m is finer than SRTM's native $30$\,m
because BlenderGIS resamples the elevation product onto its own grid. The
mesh is therefore smoother than the source data, not more accurate than it:
SRTM's specified vertical accuracy is $\pm 16$\,m, and no processing step
improves on that. What the terrain represents faithfully is the \emph{shape}
of the site --- where the ground rises and falls --- rather than absolute
height at any single point. That is what the experiment requires, since the
aircraft flies at a commanded height above its launch point and it is the
variation in ground clearance across a patrol that changes the apparent size
of a person in the image.

\paragraph{Export format.}
The terrain was exported in COLLADA rather than STL. The distinction
resolves two difficulties that arise with STL in this workflow. First, STL
carries no unit declaration, so a mesh must be rescaled by hand after import
and an error there propagates into every distance the simulation reports;
the COLLADA file declares \texttt{<unit name="meter" meter="1"/>} with a
$Z$-up axis, so Gazebo reads it at true scale unmodified. Second, STL stores
only geometry and has no texture coordinates, so an STL terrain cannot carry
aerial imagery and renders as a bare surface. COLLADA stores per-vertex
texture coordinates alongside the geometry, which is what allows the imagery
described next to be applied at all.

\paragraph{Imagery.}
Aerial imagery was exported from Google Satellite through QGIS XYZ tile
layers in Web Mercator projection, at $1792 \times 1280$ pixels, and applied
to the terrain as a texture. A cloud-free export was selected: cloud in the
imagery would appear as permanent white patches on the ground and would be a
standing source of false detections. Beyond appearance, the texture supplies
the visual clutter of vegetation, paths and roofs against which the detector
must distinguish a person --- a uniformly coloured surface would make the
detection task easier than it is in reality, and any recall figure measured
against it would be optimistic.

The imagery was initially referenced in GeoTIFF form, as exported. That file
carries EXIF metadata which Gazebo's image library does not recognise,
producing roughly thirty diagnostic messages on every load; the same image in
PNG at identical resolution was substituted, removing them at no cost in
quality. The messages were harmless, but indistinguishable at a glance from
a genuine loading failure.

% ---------------------------------------------------------------------
\subsection{Structures}
\label{sec:world-buildings}

\emph{The work described in this subsection and the next was carried out by a
project partner. It is included because the environment cannot be described
without it.}

Seventeen building footprints were traced manually in QGIS over a Google
Satellite tile layer, each carrying attributes for name, height, floor count
and roof type, and exported as GeoJSON in WGS84. The footprints were imported
into Blender, extruded to height, and given terracotta hip roofs modelled
individually.

Manual tracing was necessary because building outlines for this campus are
not available at usable quality in open mapping databases. The cost of that
choice is that footprint accuracy is limited by how precisely a polygon can
be placed over satellite imagery by hand, which is one of the two
contributions to the position error reported below.

% ---------------------------------------------------------------------
\subsection{Integration and positional accuracy}
\label{sec:world-integration}

Terrain and buildings were produced independently and against different
origins --- the terrain against $6.07778^{\circ}$\,N,
$80.20722^{\circ}$\,E and the buildings against the Administration Building
entrance --- so an offset of $1666.281$\,m east and $-166.891$\,m north was
computed and applied to bring them into one coordinate system. Each building
was then adjusted vertically to sit on the terrain slope beneath it, and the
combined scene exported as a single mesh. The model is placed at
$z = -31.412$\,m so that the terrain surface coincides with the simulator's
$z = 0$ plane.

A $2000 \times 2000$\,m invisible collision plane is placed at that same
level. Without it, an aircraft leaving the $382 \times 535$\,m modelled area
would encounter no ground at all and fall indefinitely, ending the run.

Positional accuracy was checked against two references. Against the source
GeoJSON footprints, agreement was below $0.01$\,m, confirming that the
coordinate transformation itself is correct. Against building centroids
picked independently from Google Maps, the mean error over sixteen structures
was $5.64$\,m, with a maximum of $10.59$\,m and a minimum of $0.38$\,m.

The second figure is the meaningful one, and it is consistent with what the
inputs allow: SRTM's $30$\,m cells and the precision of manual tracing
together account for errors of this order. For context, consumer GPS
receivers of the kind carried by small UAVs achieve $2$--$5$\,m in practice,
so the discrepancy between the modelled campus and the real one is of the
same order as the uncertainty a real aircraft would have in its own position.

% ---------------------------------------------------------------------
\subsection{Relationship to the experiments reported in this work}
\label{sec:world-scope}

The campus environment described above is a contribution in its own right and
is intended for mission-level UAV work at this site. The hardware-in-the-loop
measurements reported in Section~\ref{sec:hitl} were \emph{not} taken in it.
They were taken in a separate world, which was already validated against the
network-namespace and lockstep configuration those experiments depend on, and
which contains five human models at fixed coordinates that make detection
recall a counted quantity.

The two are kept distinct deliberately. Changing the world and the network
architecture at the same time would have made any change in the results
impossible to attribute. Running the hardware-in-the-loop experiments in the
campus environment is the natural next step, and requires only that the
autopilot's flight-dynamics plugin be added to that world.

\begin{figure}[htbp]
\centering
\includegraphics[width=0.9\textwidth]{world_rendered.png}
\caption{The campus environment in Gazebo. Google Satellite imagery is
applied to an SRTM-derived terrain mesh, with individually modelled buildings
placed on it. Terrain and structures both carry collision geometry generated
from the meshes used for rendering, so the surface an aircraft can strike is
the surface its camera sees.}
\label{fig:world-rendered}
\end{figure}

% =====================================================================
%  IMAGE CHECKLIST -- save into docs/report/figures/
%
%  world_rendered.png   (required)
%      Gazebo client view of the campus. Two things to fix relative to the
%      screenshot taken so far:
%        1. Pull the camera back and up. The current shot is close in on a
%           single roof; a wider view showing several buildings and the
%           terrain relief demonstrates far more.
%        2. Crop out the left World/Insert/Layers panel and the bottom
%           status bar. Keep only the 3D viewport.
%
%  terrain_wireframe.png   (optional, recommended)
%      View > Wireframe in Gazebo, camera angled low across the terrain.
%      This makes the 65x65 grid and the relief visible at once, and is the
%      only image that proves the terrain is a height map rather than a
%      textured plane -- the claim in this section a reader is most likely
%      to doubt.
%
%  NOT recommended as a figure: the GPS validation table. It is a partner's
%  result; cite the summary figures in prose as done above rather than
%  reproducing the full table under your own name.
% =====================================================================
